\section{Review of Newton's laws of motion}
    
    Before delving into a review of Newton's laws of motion, we need to be clear of the definitions of some important terms used in mechanics:
    \begin{itemize}
    	\item \emph{Particle:} an idealized object that has mass and position but occupies zero volume.
    	
    	\item \emph{Event:} a specific occurrence that happens at a unique place in space and a specific instant in time.
    	
    	\item \emph{Observer:} a person or an equipment which can locate, record, measure and interpret an event.
    	
    	\item \emph{Frame of reference:} a coordinate system combined with a clock, used by an observer to measure the position and motion of particles.
    	
    	\item \emph{Rest:} a particle is said to be at rest when its position does not change relative to a chosen frame of reference over time.
    	
    	\item \emph{Motion:} the change in position or orientation of an object with respect to a chosen frame of reference over time.
    	
    	\item \emph{Inertia:} the tendency of a body to resist any change in its motion or configuration.
    	
    	\item \emph{Position:} the location of a particle in space with respect to a fixed reference point known as the origin. If $(x,y,z)$ are the coordinates of a particle in Cartesian space, then its position vector is
    	\begin{align}
    		\vb{r}&=x\vu{i}+y\vu{j}+z\vu{k} \label{eqn:posn-vec}
    	\end{align}
    	
    	\item \emph{Displacement:} the change in the position of a particle. If a particle is initially at the position $\vb{r}_i=x_i\vu{i}+y_i\vu{j}+z_i\vu{k}$ and then finally moves to the position $\vb{r}_f=x_f\vu{i}+y_f\vu{j}+z_f\vu{k}$, its displacement vector is
    	\begin{align}
    		\Delta\vb{r}&=\vb{r}_f-\vb{r}_i=(x_f-x_i)\vu{i}+(y_f-y_i)\vu{j}+(z_f-z_i)\vu{k} \label{eqn:displ-vec}
    	\end{align}
    	If the displacement is infinitessimally small, then we can write
    	\begin{align}
    		\dd{\vb{r}}&=\vu{i}\dd{x}+\vu{j}\dd{y}+\vu{k}\dd{z} \label{eqn:displ-vec-inf}
    	\end{align}
    	
    	\item \emph{Velocity:} the rate of change of position of a particle. If the position of a particle varies with time, then its instantaneous velocity is the first derivative of its position with respect to time.
    	\begin{align}
    		\vb{v}=\dot{\vb{r}}&=\dv{\vb{r}}{t} \label{eqn:vel-vec}
    	\end{align}
    	If the positions of a particle at two instants of time are $\vb{r}(t_1)=\vb{r}_1$ and $\vb{r}(t_2)=\vb{r}_2$, then its displacement is obtained by the integral
    	\begin{align}
    		\int_{\vb{r}_1}^{\vb{r}_2}{\dd{\vb{r}}}&=\int_{t_1}^{t_2}{\vb{v}\dd{t}} \nonumber\\
    		\implies\vb{r}_2-\vb{r}_1&=\int_{t_1}^{t_2}{\vb{v}\dd{t}} \label{eqn:disp-integ}
    	\end{align}
    	
    	\item \emph{Acceleration:} the rate of change of velocity of a particle. If the velocity of a particle varies with time, then its instantaneous acceleration is the first derivative of its velocity with respect to time, or the second derivative of its position with respect to time.
    	\begin{align}
    		\vb{a}=\dot{\vb{v}}&=\dv{\vb{v}}{t} \label{eqn:acc-vec:v}\\
    		\qq{or}\vb{a}=\ddot{\vb{r}}&=\dv[2]{\vb{r}}{t} \label{eqn:acc-vec:r}
    	\end{align}
    	If the velocities of a particle at two instants of time are $\vb{v}(t_1)=\vb{v}_1$ and $\vb{v}(t_2)=\vb{v}_2$, then the change in its velocity is obtained by the integral
    	\begin{align}
    		\int_{\vb{v}_1}^{\vb{v}_2}{\dd{\vb{v}}}&=\int_{t_1}^{t_2}{\vb{a}\dd{t}} \nonumber\\
    		\implies\vb{v}_2-\vb{v}_1&=\int_{t_1}^{t_2}{\vb{a}\dd{t}} \label{eqn:vel-integ}
    	\end{align}
    \end{itemize}
    
    \subsection{Kinematics relations}
    	The kinematic relations are a set of formulae in mechanics that describe the motion of objects moving with constant acceleration, without considering the forces that cause the motion.
    	
    	\subsubsection{Velocity-time relation}
    		Consider the motion of a particle moving with a constant acceleration $\vb{a}$. Let its initial velocity be $\vb{v}_1=\vb{v}_0$ at time $t_1=0$, and its final velocity be $\vb{v}_2=\vb{v}(t)$ at some arbitrary time $t$. Then from equation (\ref{eqn:vel-integ}), we have
    		\begin{align}
    			\vb{v}(t)-\vb{v}_0&=\vb{a}\int_{0}^{t}{\dd{t}} \nonumber\\
    			\implies\vb{v}(t)-\vb{v}_0&=\vb{a}\qty(t-0) \nonumber\\
    			\implies\Aboxed{\vb{v}(t)&=\vb{v}_0+\vb{a}t} \label{eqn:kine-rel-01}
    		\end{align}
    		which is the \emph{velocity-time relation}, or the \emph{first kinematic relation}.
    		
    	\subsubsection{Position-time relation}
    		Consider the motion of a particle moving with a constant acceleration $\vb{a}$. Let its initial position be $\vb{r}_1=\vb{r}_0$ at time $t_1=0$, and its final position be $\vb{r}_2=\vb{r}(t)$ at some arbitrary time $t$. Then from equation (\ref{eqn:disp-integ}), we have
    		\begin{align}
    			\vb{r}(t)-\vb{r}_0&=\int_{0}^{t}{\vb{v}\dd{t}} \nonumber\\
    				&=\int_{0}^{t}{\qty(\vb{v}_0+\vb{a}t)\dd{t}}\qq{[using eqn. (\ref{eqn:kine-rel-01})]} \nonumber\\
    				&=\vb{v}_0\int_{0}^{t}{\dd{t}}+\vb{a}\int_{0}^{t}{t\dd{t}} \nonumber\\
    				&=\vb{v}_0\qty(t-0)+\vb{a}\left. \dfrac{t^2}{2} \right|_{0}^{t} \nonumber\\
    			\implies\Aboxed{\vb{r}(t)&=\vb{r}_0+\vb{v}_0t+\dfrac{1}{2}\vb{a}t^2} \label{eqn:kine-rel-02}
    		\end{align}
    		which is the \emph{position-time relation}, or the \emph{second kinematic relation}.
    		
    	\subsubsection{Position-velocity relation}
    		We know that
    		\begin{align*}
    			\vb{a}&=\dv{\vb{v}}{t}
    		\end{align*}
    		Taking the dot product of the above with an infinitessimal displacement $\dd{\vb{r}}$, we have
    		\begin{align}
    			\vb{a}\cdot\dd{\vb{r}}&=\dv{\vb{v}}{t}\cdot\dd{\vb{r}} \nonumber\\
    			\implies\vb{a}\cdot\dd{\vb{r}}&=\dv{\vb{r}}{t}\cdot\dd{\vb{v}} \nonumber\\
    			\implies\vb{a}\cdot\dd{\vb{r}}&=\vb{v}\cdot\dd{\vb{v}} \label{eqn:adr-vdv-01}
    		\end{align}
    		Now, using the product rule for differentiation, we have
    		\begin{align*}
    			\dd{\qty(\vb{v}\cdot\vb{v})}&=\dd{\vb{v}}\cdot\vb{v}+\vb{v}\cdot\dd{\vb{v}} \\
    			\implies 2\vb{v}\cdot\dd{\vb{v}}&=\dd{\qty(\vb{v}\cdot\vb{v})} \\
    			\implies\vb{v}\cdot\dd{\vb{v}}&=\dfrac{1}{2}\dd{\qty(\vb{v}\cdot\vb{v})}
    		\end{align*}
    		Then equation (\ref{eqn:adr-vdv-01}) becomes
    		\begin{align}
    			\vb{a}\cdot\dd{\vb{r}}&=\dfrac{1}{2}\dd{\qty(\vb{v}\cdot\vb{v})} \label{eqn:adr-vdv-02}
    		\end{align}
    		Suppose we have a particle moving under constant acceleration $\vb{a}$. If at its initial and final positions $\vb{r}_0$ and $\vb{r}$, it has initial and final velocities $\vb{v}_0$ and $\vb{v}$, then we can integrate equation (eqn:adr-vdv-02) as
    		\begin{align}
    			\vb{a}\cdot\int_{\vb{r}_0}^{\vb{r}}{\dd{\vb{r}}}&=\dfrac{1}{2}\int_{\vb{v}_0}^{\vb{v}}{\dd{\qty(\vb{v}\cdot\vb{v})}} \nonumber\\
    			\implies\vb{a}\cdot\left. \vb{r} \right|_{\vb{r}_0}^{\vb{r}}&=\dfrac{1}{2}\left. \vb{v}\cdot\vb{v} \right|_{\vb{v}_0}^{\vb{v}} \nonumber\\
    			\implies\vb{a}\cdot\qty(\vb{r}-\vb{r}_0)&=\dfrac{1}{2}\qty(\vb{v}\cdot\vb{v}-\vb{v}_0\cdot\vb{v}_0) \nonumber\\
    				&=\dfrac{1}{2}\qty(v^2-v_0^2) \nonumber\\
    			\implies\Aboxed{v^2&=v_0^2+2\vb{a}\cdot\qty(v^2-v_0^2)} \label{eqn:kine-rel-03}
    		\end{align}
    		which is the \emph{position-velocity relation}, or the \emph{third kinematic relation}.
    
    \subsection{Newton's laws of motion}
    
	    \subsubsection{Newton's first law of motion}
	    	\begin{tcolorbox}[colframe=black, colback=lightgray, opacityfill=0.2, arc=0mm]
	    		A particle tends to maintain its state of rest or uniform motion unless it is acted upon by an external agency.
	    	\end{tcolorbox}\vspace{-2.5ex}
	    	\begin{itemize}
	    		\item The importance of Newton's first law of motion lies in that it \ul{defines force} as:
	    		
	    		\emph{Force is any external agency that brings about a change in the state of rest or uniform motion of a particle.}
	    	\end{itemize}
	    
	    \subsubsection{Newton's second law of motion}
	    	\begin{tcolorbox}[colframe=black, colback=lightgray, opacityfill=0.2, arc=0mm]
	    		The rate of change of momentum of a particle is directly proportional to the force exerted on it, and takes place in the direction of the force.
	    		\begin{align*}
	    			\vb{F}&\propto\dv{\vb{p}}{t}
	    		\end{align*}
	    	\end{tcolorbox}%\vspace{-2.5ex}
	    	
	    	This implies that we have
	    	\begin{align*}
	    		\vb{F}&=k\dv{\vb{p}}{t}
	    	\end{align*}
	    	where $k$ is the proportionality constant.
	    	
	    	Now, for a particle of mass $m$ and moving with a velocity $\vb{v}$, its momentum is defined as $\vb{p}=m\vb{v}$. Hence, Newton's second law becomes
	    	\begin{align*}
	    		\vb{F}&=k\dv{t}\qty(m\vb{v}) \\
	    			&=km\dv{\vb{v}}{t}\qq{[assuming $m$ to be constant]} \\
	    		\vb{F}&=km\vb{a}
	    	\end{align*}
	    	The constant $k$ is set to $1$ by our choice of the units of force. We define $1$ newton to be that amount of force that brings about an acceleration of $1$ \unit{\metre\per{\second\squared}} on a particle of mass $1$ \unit{\kilo\gram}. Then from the above, we have
	    	\begin{align*}
	    		\qty(1\;\unit{\newton})&=k\times\qty(1\;\unit{\kilo\gram})\times\qty(1\;\unit{\metre\per{\second\squared}}) \\
	    		\implies k&=1
	    	\end{align*}
	    	Therefore, for a particle \emph{whose mass remains unchanged}, Newton's second law of motion can be written as
	    	\begin{align*}
	    		\Aboxed{\vb{F}&=m\vb{a}}
	    	\end{align*}
	    	
	    	\begin{itemize}
	    	\item The importance of Newton's second law of motion lies in that it \ul{measures force} as:
	    	
	    	\emph{If a force of magnitude $F$ brings about an acceleration $a$ on a particle of mass $m$, then the force is measured to be $F=ma$.}
	    	\end{itemize}
	    
	    \subsubsection{Newton's third law of motion}
	    	\begin{tcolorbox}[colframe=black, colback=lightgray, opacityfill=0.2, arc=0mm]
	    		For every action, there is an equal and opposite reaction, both of which act on different bodies.
	    		\begin{align*}
	    			\vb{F}_{12}&=-\vb{F}_{21}
	    		\end{align*}
	    	\end{tcolorbox}\vspace{-2.5ex}
	    	\begin{itemize}
	    	\item The importance of Newton's second law of motion lies in that it \ul{gives the origin of force} as:
	    	
	    	\emph{If a particle experiences a force of magnitude $F$, it is because it exerts another force of the same magnitude $F$ but in the opposite direction on some other particle.}
	    	\end{itemize}
	    	
	\subsection{Special cases of Newton's second law of motion}
		
		\subsubsection{The first law as a special case}
		
		\subsubsection{The third law as a special case}